% mage-auto-heading: was-implementation-ever-the-bottleneck
% mage-auto-heading: agents-entered-an-implementation-centered-discipline
% mage-auto-heading: the-old-dream-with-a-new-substrate
% mage-auto-heading: where-engineering-effort-moves
% mage-auto-heading: engineering-the-production-system
% Authoritative LaTeX manuscript; Markdown is a generated compatibility projection.

% <!-- part-title: The Profession -->

% <!-- chapter-title: The Reorganization of Software Work -->

\chapter{The Reorganization of Software Work}

\textbf{In software, MAGE is software engineering under changed economics.} The analysis in \magevalue{chapter:theory} leaves a practical question: how should the profession reorganize its work when implementation and routine analysis become cheaper? Answering it requires looking beyond the time engineers spend writing code to the larger process in which implementation, design, verification, and operation interact.

\section{Was Implementation Ever the Bottleneck?}\label{heading:was-implementation-ever-the-bottleneck}

A reasonable objection begins with a fact: \textbf{software engineers do not spend most of their workweek typing code.} This result predates generative AI. Meyer et al., drawing on 5,971 responses from professional developers~\citep{meyer2021goodday}, found that software work spans a much broader mixture of development and non-development activities. More recently, Kumar et al. separated the developer workweek into sixteen activities; in their Microsoft sample, coding new features occupied roughly 11 percent of the actual week, alongside debugging, architecture and design, code review, testing, refactoring, environment work, security, communication, and other activities~\citep{kumar2025timewarp}.

Butler et al. turn that observation into a critique of contemporary GenAI~\citep{butler2026eightmyths}. If the technology merely makes the coding slice faster, its whole-job effect has a low ceiling. Even if coding occupies only roughly 10--15 percent of the workweek, doubling performance on it cannot double developer productivity; faster source production can simply move pressure downstream into review, testing, integration, and other work. Their criticism applies well to \textbf{GenAI used as individual code-generation assistance}. It also supports a broader point MAGE shares: lines of generated code are not an engineering outcome, and handing developers a tool does not by itself redesign the system of work around them.

The disagreement begins when that fixed slice is treated as a bound on \textbf{agentic engineering}. Time categories are not independent production stages. Software engineers learn about a design by realizing it: implementation exposes missing abstractions; debugging changes what we believe about behavior; tests expose underspecified obligations; refactoring discovers structure; deployment reveals couplings that static design did not anticipate. This interleaving is not a new observation. Ralph\textquotesingle s Sensemaking--Coevolution--Implementation theory models software development as repeated movement among understanding the context, revising the problem and design together, and constructing, debugging, and deploying the artifact; a later multi-method study using more than 1,300 developers and four longitudinal cases strongly supported that process account~\citep{ralph2015sci, ralph2016processtheory}.

% <!-- box-family: inset -->
% <!-- inset-domain: SOFTWARE ENGINEERING RESEARCH -->
% mage-box-prefix: Inset —
\begin{bookbox}{inset}{Measuring Software Engineering Work}
\label{inset-measuring-software-engineering-work}
Studies of developer time describe how work is allocated, not necessarily what constrains its productivity. Kumar et al. distinguish sixteen activities in a Microsoft sample, reporting approximately 12 percent of the workweek in communication and meetings, 11 percent in coding, 9 percent in debugging, 6 percent in architecture and design, and 5 percent in pull-request and code review~\citep{kumar2025timewarp}. These categories are useful descriptions of work, but their proportions do not establish how much of the surrounding activity could change if implementation became cheaper.

The distinction matters when interpreting GenAI productivity claims. Butler et al.\textquotesingle s fixed-slice objection concerns AI used primarily to accelerate code creation while the surrounding development process remains unchanged~\citep{butler2026eightmyths}. Ralph\textquotesingle s Sensemaking--Coevolution--Implementation theory offers a different account of the process: understanding, designing, implementing, and evaluating interact rather than proceeding as independent stages~\citep{ralph2015sci}. Its empirical follow-up combined a survey of more than 1,300 developers with four longitudinal case studies~\citep{ralph2016processtheory}. Reeves made an earlier, stronger argument that implementation itself constitutes software design~\citep{reeves1992whatissoftwaredesign}.

The fraction of time spent coding is therefore not a measure of how strongly implementation constrains the rest of engineering work.
\end{bookbox}

MAGE therefore uses \textbf{implementation capacity} in an economic rather than time-sheet sense. A cost is imposed by every design hypothesis that has to be realized before the team can learn from it, even if the eventual keystrokes occupy only a small fraction of the calendar. Reduce that cost sufficiently and more alternatives become affordable, feedback arrives sooner, and intelligence can participate in more of the surrounding loop. Modern coding agents already reach beyond source generation into repository inspection, debugging, test construction, refactoring, tool execution, and multi-step changes. Whether they can safely take on still more of that loop is exactly the engineering question MAGE addresses.

This makes the disagreement empirically testable. \textbf{If assistance accelerates only code writing while the surrounding work stays unchanged, the fixed-slice argument largely holds. If governed agents can assume meaningful portions of debugging, testing, refactoring, integration, migration, and operation without proportionate growth in human review, a fixed coding-time slice is no longer the right model of attainable leverage.} The relevant boundary is the work delegated and the supervision it requires. MAGE asks how to cross that boundary without losing engineering control.

\section{Agents Entered an Implementation-Centered Discipline}\label{heading:agents-entered-an-implementation-centered-discipline}

% <!-- point: agents-inherited-the-implementation-centered-object | The first agentic-SE research inherited the engineering object in front of it — a repository and a task — and asked how to make a stronger reasoner perform it. -->

The first generation of agentic software-engineering research inherited the engineering object already in front of it: a repository and a software-development task. The central question was how to make an increasingly capable reasoner perform that task more effectively.

% <!-- point: repocoder-sweagent-autocoderover-share-an-orientation | Retrieval, purpose-built interfaces, program analysis, and specification inference differ, but they share an orientation: improve the machinery an agent uses to act on the software artifact. -->

The progression is instructive. Early agentic software-engineering research improved agents\textquotesingle{} ability to work with existing software artifacts. Repository retrieval made relevant code and documentation available during implementation~\citep{zhang2023repocoder}. Purpose-built interfaces improved an agent\textquotesingle s ability to navigate repositories, edit files, and execute tests~\citep{yang2024sweagent}. Program analysis helped agents locate and reason about changes by exploiting program structure rather than treating a project as a collection of files~\citep{zhang2024autocoderover}. More recent work has also attempted to infer specifications from existing artifacts~\citep{ruan2025specrover}. These developments improve how agents understand and modify software, and establish important foundations for autonomous implementation. They share an orientation: improving the machinery an agent uses to act on the software artifact.

% <!-- point: agent-governance-research-is-another-antecedent | Agent-governance research translates obligations into controls across design-time, runtime, and assurance layers, evaluates actions against path-dependent policies, and integrates enforcement into an agent's operating environment; MAGE's governance-conversion process adds how implementation reveals missing obligations. | terms: governance-conversion, alignment-principle -->

Research on agent governance provides another important antecedent. Existing methods translate organizational and regulatory obligations into technical controls and assign those controls across design-time, runtime, and assurance layers~\citep{koch2026governance}. Other frameworks evaluate a proposed action against policies that depend on the agent\textquotesingle s execution history~\citep{kaptein2026runtime}. Governance-by-construction systems demonstrate how policy enforcement can be integrated into an agent\textquotesingle s operating environment~\citep{shlomov2026construction}. These approaches develop architectures and methods for governing agentic systems. MAGE\textquotesingle s governance-conversion process additionally accounts for how failures encountered during implementation reveal missing obligations or controls, and how engineers preserve those lessons for subsequent work.

% <!-- point: mage-accepts-machinery-and-changes-the-object | MAGE accepts that machinery and enlarges the engineering object: what durable environment should be engineered around autonomous implementation. -->

MAGE accepts that machinery and enlarges the engineering object. The question is no longer only how an agent should navigate a repository, retrieve context, use tools, or construct a patch. It is what durable environment should be engineered around autonomous implementation: which system properties deserve explicit representation; which obligations deserve enforcement; what evidence work must produce; how representations remain in correspondence with the implementation; and how recurring judgment becomes machinery inherited by later work.

% <!-- point: the-difference-is-level-not-opposition | Retrieval, code search, interfaces, tests, and tools belong inside a governed engineering environment — pieces of a larger engineered substrate, not accessories to a coding agent. -->

The difference is one of level, not opposition. Repository retrieval, structured code search, agent-computer interfaces, tests, and tool use all belong inside a governed engineering environment. MAGE\textquotesingle s claim is that they are pieces of a larger engineered substrate rather than accessories to a coding agent.

\section{The Old Dream, with a New Substrate}\label{heading:the-old-dream-with-a-new-substrate}

% <!-- point: the-old-dream-with-a-new-substrate | The intelligent-programming-partner dream is decades old; commodity intelligence finally lets us run much of it at production scale, where its next-order problems become visible. -->

The ambition predates today\textquotesingle s coding agents. In the 1980s, Charles Rich and Richard Waters\textquotesingle s Programmer\textquotesingle s Apprentice envisioned an AI system not merely as a programming tool but as an agent in the software process. The engineer would share programming knowledge with the Apprentice and delegate work across implementation, design, requirements, verification, explanation, and documentation~\citep{rich-waters}. Its internal representations made programming concepts explicit so that routine questions could be answered from represented knowledge rather than reconstructed from first principles.

The project realized important pieces of that vision, but under different constraints. Its intelligence depended heavily on programming knowledge deliberately encoded into specialized representations; Rich and Waters themselves observed that the economics failed when too much knowledge had to be supplied for each problem. Commodity intelligence changes that premise. A general-purpose reasoner can now arrive with substantial prior competence, enter an unfamiliar repository, use tools, interpret natural-language intent, and construct substantial implementations without the engineering organization first encoding everything it might need to know.

Earlier systems could demonstrate pieces of that vision; today\textquotesingle s systems realize enough of it to expose the next-order engineering problems: keeping representations trustworthy, deciding which obligations deserve enforcement, producing evidence at agentic velocity, and converting recurring judgment into machinery that later work inherits.

\section{Where Engineering Effort Moves}\label{heading:where-engineering-effort-moves}

% <!-- point: allocation-changes-even-where-the-eventual-organization-is-uncertain | As implementation becomes cheaper the allocation of engineering work changes, while the profession’s eventual organization remains uncertain. Human effort may move toward intent, abstraction, evidence, tradeoffs, and acceptance. -->

% <!-- index-def: engineers-seat -->

As commodity intelligence assumes more implementation and routine engineering work, human effort shifts toward directing that work and governing its results. \autoref{where-effort-moves} illustrates two related changes: delegation transfers work to agents, while governance conversion preserves recurring engineering judgment in durable structure. The latter calls for Modeling to represent consequential knowledge and obligations, and Alignment to establish evidence, constraints, and controls. These are broad directions rather than a fixed allocation of effort; their relative importance depends on which capabilities improve and what investments become economical.

\begin{figure}[htbp]
\centering
\includegraphics[width=\linewidth]{book/latex/assets/where-effort-moves-ff35e2dbed.pdf}
\caption{\emph{Where Engineering Effort Moves.} As the fleet assumes work across the lifecycle, the engineer increasingly works through the surrounding engineering environment: specifying intent, choosing consequential abstractions, evaluating evidence, resolving tradeoffs, and adapting the environment itself.}
\label{where-effort-moves}
\end{figure}

% <!-- point: where-displaced-effort-settles-depends-on-what-becomes-cheaper | Where displaced effort settles depends on which parts of the production system become cheaper and more reliable, and the theory of Chapter 6 names the quantities those improvements act on. -->

Where that displaced effort settles depends on which parts of the production system become cheaper and more reliable. \magevalue{chapter:theory} supplies the distinctions. Faithful production through a representation \(R\) depends on whether it encodes the consequential obligations (\(E\)), expressed as \(P(E \mid R)\), and whether the agent fulfills them (\(F\)), expressed as \(P(F \mid E,R)\). Interpretation contributes to the second term. Governance separately concerns which unacceptable results the mechanisms detect, reject, or prevent. Cost includes realization attempts and the work required to keep representations and controls useful as the software changes. The engineer chooses the representation and the Alignment strategy together against those quantities.

% <!-- point: four-trajectories-are-the-independent-variables-behind-the-uncertainty | Technological progress can move those terms independently: reasoners, representation machinery, realization machinery, and governance machinery can each improve at their own rate, and they are the independent variables that prevent a single confident forecast. -->

Technological progress can move these terms independently. Reasoners can improve. So can the machinery for constructing and maintaining representations, the machinery for realizing them, and the machinery for detecting and refusing unacceptable results. These four trajectories can advance independently. \autoref{effort-trajectories} sketches how each would redirect engineering effort. A technical advance could change several terms at once, and representation machinery changes probabilities and costs alike.

{\footnotesize
\begin{longtable}[]{@{}
  >{\raggedright\arraybackslash}p{(\linewidth - 4\tabcolsep) * \real{0.3333}}
  >{\raggedright\arraybackslash}p{(\linewidth - 4\tabcolsep) * \real{0.3333}}
  >{\raggedright\arraybackslash}p{(\linewidth - 4\tabcolsep) * \real{0.3333}}@{}}
\caption[How technological progress can redirect engineering effort]{\emph{Possible movements of engineering effort.} Improvements in different parts of the production system act on different terms of the \magevalue{chapter:theory} engineering problem. The trajectories can proceed independently and in combination, so the rows describe the pressures each would create rather than a forecast of a single future process.}\label{effort-trajectories}\tabularnewline
\toprule\noalign{}
\begin{minipage}[b]{\linewidth}\raggedright
Technological trajectory and terms affected
\end{minipage} & \begin{minipage}[b]{\linewidth}\raggedright
What becomes economical
\end{minipage} & \begin{minipage}[b]{\linewidth}\raggedright
Engineering effort moves toward
\end{minipage} \\
\midrule\noalign{}
\endfirsthead
\toprule\noalign{}
\begin{minipage}[b]{\linewidth}\raggedright
Technological trajectory and terms affected
\end{minipage} & \begin{minipage}[b]{\linewidth}\raggedright
What becomes economical
\end{minipage} & \begin{minipage}[b]{\linewidth}\raggedright
Engineering effort moves toward
\end{minipage} \\
\midrule\noalign{}
\endhead
\bottomrule\noalign{}
\endlastfoot
Reasoners improve: \(P(F \mid E,R)\uparrow\) & Reasoning successfully from less complete representations, and leaving more legitimate choices unresolved until realization & Choosing consequential distinctions, bounding degrees of freedom, and judging results rather than specifying every realization choice \\
Representation machinery improves: \(P(E \mid R)\uparrow\), \(P(F \mid E,R)\uparrow\), lower upkeep & Constructing, maintaining, reconciling, and querying explicit representations & Choosing abstractions, validating representations, and maintaining correspondence among them and with the system \\
Realization machinery improves: \(P(F \mid E,R)\uparrow\), lower realization cost & Reliably transforming sufficiently complete representations into implementations, potentially deterministically & Specifying and validating what should be realized rather than supervising how it is implemented \\
Governance machinery improves: less escaped failure, lower upkeep & Detecting, rejecting, or preventing unacceptable realizations independently of the producing reasoner & Choosing which obligations require enforcement, placing controls, and judging residual risk \\
\end{longtable}
}

These trajectories can proceed independently or in combination. Each changes where engineering effort is most valuable:

% <!-- point: better-reasoners-move-effort-toward-choosing-distinctions-not-away-from-encoding-them | Better reasoners leave more legitimate freedom unresolved; even perfect realization does not establish that the encoding captures the right obligations, so capability does not retire engineering responsibility. -->
% <!-- point: cheap-representation-machinery-makes-externalizing-consequential-knowledge-economical | If explicit representations become cheap to construct, maintain, and reconcile, engineers can profitably externalize more consequential knowledge — choosing, creating, maintaining, and connecting the representations through which humans and machines reason about the system. -->
% <!-- point: representation-engineering-broader-than-context-engineering | Retrieval selects information for a reasoner; representation engineering chooses the form in which the engineering question is expressed. The distinction concerns access and representation, which can be combined. -->
% <!-- point: cheap-reliable-realization-moves-effort-upstream | If realization from a sufficiently complete representation becomes cheap and reliable, engineering effort moves further upstream, toward specifying and validating what should be realized rather than supervising how it is implemented. -->
% <!-- point: cabot-is-one-plausible-equilibrium-not-the-prediction | Model-driven engineering is one plausible destination, reached when representation and realization machinery improve together; MAGE includes that equilibrium without predicting it, and the opposite combination would make separately maintained models less attractive. -->
% <!-- point: cheap-controls-let-engineering-tolerate-freedom-without-granting-it-authority | If strong controls become cheap, engineers can tolerate more freedom in realization without granting that freedom equal authority over consequences; effort then moves toward choosing which obligations deserve enforcement, placing the controls, and judging the residual risk. -->

\begin{enumerate}
\item
  \textbf{Better reasoners.} As agents become more capable, engineers can leave more legitimate realization choices unresolved. The engineer states the consequential distinctions and lets the agent choose among acceptable implementations. As \magevalue{chapter:theory} showed, when conditional realization approaches certainty, faithful production approaches the probability that the representation encoded the right obligations. Better agents may infer omitted intent, but reliable realization does not establish that the obligations were right. Engineering effort therefore moves toward deciding which distinctions matter and which freedoms to preserve.
\item
  \textbf{Better representation machinery.} Cheaper construction, maintenance, and reconciliation make explicit representations economical for more engineering questions. Retrieval and representation engineering remain different: retrieval selects information for a reasoner, while representation engineering changes the form in which the question is expressed. An agent may reconstruct dependencies from source or work over a dependency model that makes those relationships explicit. Program comprehension and architecture recovery offer established techniques for obtaining such views~\citep{chikofsky1990taxonomy, ducasse2009architecture, cornelissen2009dynamic}. Engineers must decide whether to maintain a representation, derive it when needed, or reason directly over implementation. Recovery can improve realization without establishing that the recovered behavior is what the system ought to do.
\item
  \textbf{Better realization machinery.} When a sufficiently complete representation can be transformed into an implementation cheaply and reliably, engineering effort moves toward specifying and validating that representation. Deterministic generation becomes attractive where little consequential freedom remains. Model-driven engineering is one possible destination. Cabot\textquotesingle s vibe-driven model-based engineering proposes a contemporary route in which AI helps construct models and deterministic generators realize them~\citep{cabot2026vibe}. MAGE includes this arrangement without requiring it.
\item
  \textbf{Better governance machinery.} Cheaper, more reliable validators, monitors, and admission checks make independent enforcement economical for more obligations. Engineers can permit greater freedom in realization while requiring evidence that consequential properties hold. The engineering work moves toward choosing which obligations deserve enforcement, placing the mechanisms that govern them, and judging residual risk. Stronger controls do not automatically imply more model-driven development; an organization can become more governed without making models the primary source of implementation.
\end{enumerate}

% <!-- point: combinations-land-differently-for-different-properties-of-one-system | Real environments combine the trajectories, often differently for different properties of the same system, so the unit of analysis is usually the engineering question or obligation rather than the development methodology. -->
Real engineering environments will combine these trajectories, often differently for different properties of the same system. A system might maintain an authoritative state model where invalid transitions are consequential, derive architectural views when needed, allow agents substantial implementation freedom, and enforce security and dependency obligations at admission. The appropriate unit of analysis is therefore usually the engineering question or obligation, not the development methodology.

% <!-- point: more-modeling-need-not-mean-more-permanently-maintained-models | More modeling need not mean more permanently maintained models: construction cost can be repaid through repeated reuse or one consequential answer; cheap machine reasoning expands the one-off questions for which a purposeful reduction pays. -->
More Modeling need not mean more permanently maintained models. Agents may construct a dependency graph, state abstraction, or quantitative model for one consequential question and discard it afterward. A representation can repay its construction cost through repeated reuse or by answering a sufficiently important question once.

% <!-- point: recurring-judgment-becomes-engineering-capital-under-any-trajectory | Recurring judgment remains an opportunity for governance conversion; conversion and ratcheteering let later work inherit it while keeping the resulting structure open to correction. -->
Whichever trajectory dominates, recurring judgment remains an opportunity for governance conversion. The question is whether later work should inherit that judgment through the engineering environment or whether it should be supplied again by a person or agent. The answer is not always mechanization: guidance, expert review, and probabilistic evaluation remain appropriate where the obligation or its economics warrant them.

\section{Engineering the Production System}\label{heading:engineering-the-production-system}

% <!-- point: two-places-engineering-attention-can-go | Chapter 5 distinguished two places engineering attention can go as realization becomes abundant: operating agentic machinery, and engineering the production system around it so selected knowledge, evidence, and controls persist across acts of realization. -->

\magevalue{chapter:software-factories} distinguished two places engineering attention can go as realization becomes abundant. Engineers can spend increasing effort operating agentic machinery: supplying work and context, coordinating agents, inspecting output, and resolving exceptions. They can also engineer the production system around that machinery so that selected knowledge, evidence, and controls persist across acts of realization.

DevOps offers a particularly close precedent. Infrastructure as Code and Policy as Code moved knowledge that once lived in operating practice into executable artifacts. MAGE generalizes the move without insisting on code as the form: an engineering decision might become a type, model, schema, validator, permission, workflow, or documented review obligation. Useful engineering knowledge increasingly becomes structure later work can inherit; residual semantic judgment remains with people.

Code review provides one example of how this changes engineering practice. At low change volume, an experienced engineer can treat the diff as the principal reasoning surface: inspect the implementation, reconstruct the surrounding architecture, examine the tests, and decide whether the change belongs. That strategy becomes increasingly costly as implementation grows abundant. Earlier empirical work found that larger changes taxed human review and reduced useful feedback~\citep{bosu2015useful}. Greiler\textquotesingle s SCOPE proposal responds to agentic development by moving detailed implementation review into the developer--agent loop and making team review risk-proportional~\citep{greiler2026scope}. MAGE additionally changes the surfaces available to the reviewer. When the environment holds explicit models of structure, behavior, ownership, execution, measurement, and provenance, those models become review surfaces in their own right.

The reviewer can then concentrate on four questions:

\begin{itemize}
\tightlist
\item
  \textbf{Is the intended change desirable?}
\item
  \textbf{Do the models capture the relevant concern?}
\item
  \textbf{Is the evidence sufficient?} Do the checks that ran actually bound the risk this change carries?
\item
  \textbf{Should the tradeoff be accepted?}
\end{itemize}

Implementation becomes one evidentiary surface among several, rather than the origin from which every system property must be reconstructed. The reviewer still examines implementation details when they are necessary to evaluate the risk. The models supply an index into the engineering significance of a change. In the terms of \magevalue{chapter:theory}, these representations can reduce the cost of evaluating a realization without necessarily improving the probability that the agent produces an acceptable one.

% <!-- point: not-a-new-lifecycle-phase-but-a-new-means-of-exercising-it | This is not a new lifecycle phase: requirements, design, validation, operation, and maintenance remain, but engineers increasingly exercise those responsibilities through the production system in which agents work. -->

This is not a new lifecycle phase. Requirements, design, validation, operation, and maintenance remain. What changes is that engineers increasingly exercise those responsibilities through the production system in which agents work. The next section places that arrangement within the longer history of engineering through explicit representations and controls.
